
\textbf{II. PHẦN TỰ LUẬN}
\begin{ex}
	Giải hệ phương trình:
	$$ \begin{cases} 2x+y=9 \\ 3x-2y=10 \end{cases} $$
	\loigiai{
		Ta có hệ phương trình:
		$$ \begin{cases} 2x+y=9 \quad &(1) \\ 3x-2y=10 \quad &(2) \end{cases} $$
		Từ phương trình (1), suy ra $y = 9-2x$.
		Thế $y = 9-2x$ vào phương trình (2):
		$3x - 2(9-2x) = 10$
		$3x - 18 + 4x = 10$
		$7x = 10 + 18$
		$7x = 28$
		$x = \frac{28}{7}$
		$x = 4$.
		Thay $x=4$ vào $y = 9-2x$:
		$y = 9 - 2(4) = 9 - 8 = 1$.
		Vậy nghiệm của hệ phương trình là $(x; y) = (4; 1)$.
		\\ \textbf{Đáp số: $(x;y) = (4;1)$}
	}
\end{ex}

\begin{ex}
	Cho $x>0$. Rút gọn biểu thức:
	$$ P = \left(\frac{1}{\sqrt{x}+2} + \frac{1}{x+2\sqrt{x}}\right) : \frac{\sqrt{x}+1}{2\sqrt{x}+4} $$
	\loigiai{
		Điều kiện: $x>0$.
		Ta có:
		\begin{align*} P &= \left(\frac{1}{\sqrt{x}+2} + \frac{1}{\sqrt{x}(\sqrt{x}+2)}\right) : \frac{\sqrt{x}+1}{2(\sqrt{x}+2)} \\ &= \left(\frac{\sqrt{x}}{\sqrt{x}(\sqrt{x}+2)} + \frac{1}{\sqrt{x}(\sqrt{x}+2)}\right) \cdot \frac{2(\sqrt{x}+2)}{\sqrt{x}+1} \\ &= \frac{\sqrt{x}+1}{\sqrt{x}(\sqrt{x}+2)} \cdot \frac{2(\sqrt{x}+2)}{\sqrt{x}+1} \end{align*}
		Vì $x>0$ nên $\sqrt{x}+1 > 0$ và $\sqrt{x}+2 > 0$.
		\begin{align*} P &= \frac{\cancel{(\sqrt{x}+1)}}{\sqrt{x}\cancel{(\sqrt{x}+2)}} \cdot \frac{2\cancel{(\sqrt{x}+2)}}{\cancel{(\sqrt{x}+1)}} \\ &= \frac{2}{\sqrt{x}} \end{align*}
		Vậy $P = \frac{2}{\sqrt{x}}$.
		\\ \textbf{Đáp số: $P=\frac{2}{\sqrt{x}}$}
	}
\end{ex}

\begin{ex}
	Biết phương trình bậc hai $x^2-4x+1=0$ có hai nghiệm phân biệt $x_1, x_2$. Không giải phương trình. Hãy tính giá trị của biểu thức: $F=x_1^3+x_2^3+3x_1x_2(x_1+x_2+2)$.
	\loigiai{
		Phương trình $x^2-4x+1=0$ có $a=1, b=-4, c=1$.
		$\Delta' = (-2)^2 - 1 \cdot 1 = 4-1=3 > 0$, nên phương trình có hai nghiệm phân biệt $x_1, x_2$.
		Theo hệ thức Vi-ét, ta có:
		$$ \begin{cases} x_1 + x_2 = -\frac{b}{a} = -\frac{-4}{1} = 4 \\ x_1 x_2 = \frac{c}{a} = \frac{1}{1} = 1 \end{cases} $$
		Ta biến đổi biểu thức $x_1^3+x_2^3$:
		$x_1^3+x_2^3 = (x_1+x_2)^3 - 3x_1x_2(x_1+x_2)$.
		Thay vào biểu thức F:
		\begin{align*} F &= (x_1+x_2)^3 - 3x_1x_2(x_1+x_2) + 3x_1x_2(x_1+x_2+2) \\ &= (x_1+x_2)^3 - 3x_1x_2(x_1+x_2) + 3x_1x_2(x_1+x_2) + 3x_1x_2(2) \\ &= (x_1+x_2)^3 + 6x_1x_2 \end{align*}
		Thay giá trị từ hệ thức Vi-ét vào biểu thức F:
		\begin{align*} F &= (4)^3 + 6(1) \\ &= 64 + 6 \\ &= 70 \end{align*}
		Vậy $F=70$.
		\\ \textbf{Đáp số: $F=70$}
	}
\end{ex}

\begin{ex}
	Một công viên hình chữ nhật có chiều dài hơn chiều rộng 20 mét và có chu vi bằng 160m.
	\begin{enumerate}
		\item[a)] Tính chiều dài và chiều rộng của công viên?
		\item[b)] Công viên được lát bằng gạch chống trơn và chi phí lát gạch chống trơn cho 1m$^2$ là 350 nghìn đồng. Hỏi kinh phí lát gạch chống trơn của khuôn viên là bao nhiêu tiền?
	\end{enumerate}
	\loigiai{
		\begin{enumerate}
			\item[a)] Gọi chiều dài của công viên là $x$ (mét) và chiều rộng của công viên là $y$ (mét).
			Điều kiện: $x > 0, y > 0, x > y$.
			
			Theo đề bài, chiều dài hơn chiều rộng 20 mét, ta có phương trình:
			$x - y = 20 \quad (1)$
			
			Chu vi của công viên hình chữ nhật là 160m, ta có phương trình:
			$2(x+y) = 160$
			$x+y = \frac{160}{2}$
			$x+y = 80 \quad (2)$
			
			Từ (1) và (2), ta có hệ phương trình:
			$$ \begin{cases} x - y = 20 \\ x + y = 80 \end{cases} $$
			Cộng vế theo vế của hai phương trình, ta được:
			$(x-y) + (x+y) = 20 + 80$
			$2x = 100$
			$x = \frac{100}{2}$
			$x = 50$.
			Thay $x=50$ vào phương trình (2):
			$50 + y = 80$
			$y = 80 - 50$
			$y = 30$.
			Vậy, chiều dài của công viên là 50 mét và chiều rộng là 30 mét.
			(Thỏa mãn điều kiện $x>0, y>0, x>y$).
			
			\item[b)] Diện tích của công viên hình chữ nhật là:
			$S = x \times y = 50 \times 30 = 1500$ (m$^2$).
			Chi phí lát gạch chống trơn cho 1m$^2$ là 350 nghìn đồng.
			Tổng kinh phí lát gạch chống trơn cho toàn bộ khuôn viên là:
			$K = 1500 \text{ m}^2 \times 350 \text{ nghìn đồng/m}^2$
			$K = 525000$ nghìn đồng.
			$K = 525\,000\,000$ đồng.
			Vậy, kinh phí lát gạch chống trơn của khuôn viên là 525.000.000 đồng.
		\end{enumerate}
	}
\end{ex}

\begin{ex}
	Xét phép thử “Viết ngẫu nhiên một số tự nhiên lẻ có hai chữ số”. Xác suất của biến cố B: “Số tự nhiên được viết ra là bội của 9” là bao nhiêu?
	\loigiai{
		Không gian mẫu: $\Omega = \{11; 13; 15; \dots ; 97; 99\}$.\\
		Số phần tử của không gian mẫu: $n(\Omega) = \frac{99-11}{2} + 1 = 45$.\\
		Biến cố B: “Số tự nhiên được viết ra là bội của 9”.\\
		Các phần tử của biến cố B (số tự nhiên lẻ có hai chữ số là bội của 9):\\
		$B = \{27; 45; 63; 81; 99\}$.
		Số phần tử của biến cố B: $n(B) = 5$.\\
		Xác suất của biến cố B:\\
		$P(B) = \frac{n(B)}{n(\Omega)} = \frac{5}{45} = \frac{1}{9}$.
	}
\end{ex}
\begin{ex}
	Cho tam giác nhọn ABC ($AB \neq AC$). Vẽ đường tròn tâm O đường kính BC cắt hai cạnh AB và AC lần lượt tại F và E. BE cắt CF tại H.
	\begin{enumerate}
		\item[a)] Chứng minh bốn điểm A, E, H, F cùng nằm trên một đường tròn.
		\item[b)] Tia AH cắt BC tại D. Chứng minh $DH \cdot DA = DB \cdot DC$ và OE là các tiếp tuyến của đường tròn ngoại tiếp tam giác EHF?
	\end{enumerate}
	\loigiai{
	\begin{center}
		

				
		
	\end{center}
	}
\end{ex}

\begin{ex}
	Cho a, b, c là các số dương thỏa mãn $a^2+b^2+c^2=1$. Tìm giá trị lớn nhất của biểu thức $A = (1+2a)(1+2bc)$.
\end{ex}

\setcounterref{ex}{0}
